\section{Discusión}

\textbf{La brecha entre exactitud estricta y laxa} mide cuánto penaliza la
coincidencia textual exacta a respuestas que un juez semántico considera
equivalentes, y la taxonomía la descompone en dos causas: ausencia de error
semántico, y sustitución de un valor del vocabulario cerrado por un sinónimo.

\textbf{Endurecer el prompt no elimina el problema.} Aunque el prompt declara
que un sinónimo de un valor categórico cerrado es una violación de formato
(Sección \ref{sec:prompting}), \texttt{uso\_de\_sinonimos} aparece en
prácticamente todos los modelos: copiar textualmente un token de una lista
cerrada exige un tipo de obediencia distinto del que exige acertar el
contenido, y varios modelos con buena exactitud de campo prefieren la forma
natural (``tele'') sobre la literal exigida (``tv''). El caso extremo es
\texttt{SmolLM2-360M-Instruct}, cuyas respuestas incorrectas se explican casi
por completo por esta única causa (Tabla \ref{tab:taxonomia}).

\textbf{Tamaño, arquitectura y generación.} Que \texttt{Qwen3.5-0.8B} supere
a cuatro de los seis modelos del tier superior con menos de mil millones de
parámetros muestra que el desempeño no depende sólo del conteo de parámetros,
y los +25,0 puntos de \texttt{granite-4.0-1b} sobre el mejor modelo de la
generación anterior cuantifican el desplazamiento de lo factible en el rango
sub-2B.

\textbf{Penalización de latencia de la atención híbrida.}
\label{sec:discusion-latencia-hibrida}
\texttt{granite-4.0-h-350m} tarda \textbf{56,45} segundos por comando frente
a los \textbf{13,77} de la variante densa \texttt{granite-4.0-350m}
(\textbf{$\times$4,10}), y \texttt{granite-4.0-h-1b} \textbf{123,15} frente a
los \textbf{42,78} de \texttt{granite-4.0-1b} (\textbf{$\times$2,88}). Que se
reproduzca en \textbf{ambos} pares de tamaño, con la misma familia y el mismo
protocolo, indica un efecto arquitectónico y no un artefacto puntual. Las
variantes híbridas se promocionan por su complejidad sub-cuadrática en la
longitud de secuencia, beneficio que se materializa en secuencias largas;
aquí los prompts son \textbf{cortos} y la ejecución es sobre \textbf{CPU}, de
modo que ese beneficio asintótico no llega a manifestarse mientras que la
sobrecarga \textbf{constante} de gestionar dos tipos de caché —atención
estándar y estado lineal/recurrente— se paga en cada comando y domina el
tiempo total. \textbf{A igual cantidad de parámetros, la variante densa es
preferible en CPU para prompts cortos.}

\textbf{Viabilidad práctica.} Los 42,78 segundos por comando del modelo más
exacto son incompatibles con una interacción de voz en tiempo real, aunque
aceptables para automatización diferida; entre los sub-1B,
\texttt{Qwen3.5-0.8B} combina 75,0\% de exactitud estricta con 18,17
segundos, el compromiso más equilibrado del roster.
